\section*{Chapter \ref{ch:iv:probability-ii} --- Probability II}

\begin{solution}{iq:iv:probability-ii:1}
One. Order $i$ is in its original place with probability $\tfrac1{20}$, and there are twenty orders; linearity
does not care that the events are dependent. (The count is approximately Poisson with mean 1.)
\iqlookfor{indicators and linearity, without enumerating permutations.}
\end{solution}

\begin{solution}{iq:iv:probability-ii:2}
The first six takes a geometric number of rolls with mean 6. The second six needs another independent geometric
wait, so 12 in all.
\iqlookfor{the geometric mean and additivity of independent waits.}
\end{solution}

\begin{solution}{iq:iv:probability-ii:3}
A run starts at toss 1 and at each toss $i \ge 2$ that differs from toss $i-1$, which happens with probability
$\tfrac12$: $1 + 9 \times \tfrac12 = 5.5$. Enumerating all 1\,024 sequences gives the same.
\iqlookfor{counting run starts with indicators.}
\end{solution}

\begin{solution}{iq:iv:probability-ii:4}
The two points cut $[0,1]$ into three exchangeable pieces of mean $\tfrac13$, so the smaller has mean $\tfrac13$ and
the larger $\tfrac23$.
\iqlookfor{the spacings symmetry rather than an integral (which gives the same).}
\end{solution}

\begin{solution}{iq:iv:probability-ii:5}
The first of two exponential clocks rings with probability proportional to its rate: $3/(3+1) = \tfrac34$ for A.
The waiting time is exponential with rate 4 a minute: 15 seconds on average.
\iqlookfor{the competing-exponentials rule and the rate of the minimum.}
\end{solution}

\begin{solution}{iq:iv:probability-ii:6}
After the first roll, each roll matches the previous one with probability $\tfrac16$, whatever it was, so the
wait is $1 + 6 = 7$. For two sixes in a row, each attempt needs the previous roll to be a six as well, which
happens only a sixth of the time, and a failure on the second roll of an attempt sends you back to the start:
the states carry more memory, and the wait is $6 + 36 = 42$.
\iqlookfor{recognising the one-state structure, and explaining the difference with the two-state chain.}
\end{solution}

\begin{solution}{iq:iv:probability-ii:7}
With $r = q/p = 49/51$ and $a = b = 10$: $\big(1 - r^{10}\big)/\big(1 - r^{20}\big) = 1/(1 + r^{10}) \approx 1/(1 +
0.670) \approx 0.599$. A one-point edge per bet becomes a ten-point edge over the game, because the game lasts
about a hundred bets.
\iqlookfor{the ruin formula with drift, and an intuition for why small edges compound.}
\end{solution}

\begin{solution}{iq:iv:probability-ii:8}
$S_t^2 - t$ is a martingale; optional stopping at the exit gives $\E[\tau] = \E[S_\tau^2] - 9 = 100 \times
\tfrac3{10} - 9 = 21$, which is $x(n - x) = 3 \times 7$. A simulation of 20\,000 walks agrees.
\iqlookfor{the quadratic martingale, or the first-step recursion, and the product formula.}
\end{solution}

\begin{solution}{iq:iv:probability-ii:9}
Mean $40 \times 0.3 = 12$. Variance $\E[N]\Var(X) + \Var(N)\E[X]^2 = 40 \times 4 + 40 \times 0.09 = 163.6$ (for a
Poisson count $\Var(N) = \E[N]$), so a standard deviation of about 12.8 thousand. The day's Sharpe ratio is about
0.94.
\iqlookfor{Wald's identity and the compound variance, with the Poisson simplification.}
\end{solution}

\begin{solution}{iq:iv:probability-ii:10}
$365\big(1 - (364/365)^{20}\big) \approx 19.49$: on average about half a collision among twenty keys, consistent
with the birthday problem.
\iqlookfor{indicators for each bucket being used.}
\end{solution}

\begin{solution}{iq:iv:probability-ii:11}
HHT comes first with probability $\tfrac23$: once HH has appeared, HHT is certain before HTH (every T completes
HHT), whereas an HT can be spoiled by a second T. On its own, HHT takes 8 tosses on average and HTH takes 10,
since HTH overlaps itself (its final H starts a new attempt). Here both answers favour HHT; in general the race
and the waiting time can disagree (THH beats HHT three times in four, with equal waits of 8), which is the
lesson of Penney's game: the waiting time is a property of one pattern, the race a property of the pair.
\iqlookfor{first-step analysis for the race, the overlap argument for the waits, and not confusing the two.}
\end{solution}

\begin{solution}{iq:iv:probability-ii:12}
Backward induction: with one offer left its value is $V_1 = \tfrac12$; with $k$ left, accept an offer above
$V_{k-1}$, so $V_k = \E[\max(U, V_{k-1})] = (1 + V_{k-1}^2)/2$. $V_2 = 0.625$, $V_3 \approx 0.695$, $V_4
\approx 0.742$, $V_5 \approx 0.775$. Accept the first offer if it exceeds 0.742, the second if above 0.695, the
third if above 0.625, the fourth if above 0.5, and take the fifth.
\iqlookfor{backward induction on continuation values with decreasing thresholds.}
\end{solution}

\begin{solution}{iq:iv:probability-ii:13}
The second arrival has mean $2/6$ of an hour, 20 minutes. The six gaps average 10 minutes, but the gap covering a
random moment has expected length $\sum \E[S_i^2] = 6 \times \tfrac{2}{6 \times 7} = \tfrac27$ of an hour, about 17
minutes: a gap is sampled with probability proportional to its length. Monitoring that measures gaps from
random moments overstates the typical gap by this length bias. A simulation of 200\,000 hours gives 0.286.
\iqlookfor{order-statistic means, and the length-biased sampling behind the inspection paradox.}
\end{solution}

\begin{solution}{iq:iv:probability-ii:14}
Couple the coins: use the same uniform $U_i$ for toss $i$ of both, with heads when $U_i < p$. Every head of the
fair coin ($U_i < 0.5$) is a head of the biased coin ($U_i < 0.6$), so on every outcome the biased coin has at
least as many heads, and the event ``at least seven'' is at least as likely for it. The values are 0.382 and
$176/1\,024 \approx 0.172$.
\iqlookfor{constructing a monotone coupling, then confirming by the binomial sums.}
\end{solution}
